\documentclass[11pt,oneside]{book}
\usepackage[margin=1in]{geometry}
\usepackage[T1]{fontenc}
\usepackage[utf8]{inputenc}
\usepackage{lmodern}
\usepackage{microtype}
\usepackage{amsmath,amssymb,amsthm,mathtools}
\usepackage{booktabs,longtable,array,tabularx}
\usepackage{enumitem}
\usepackage{hyperref}
\usepackage{bookmark}
\usepackage{fancyhdr}
\usepackage{setspace}
\usepackage{xcolor}

\hypersetup{
  colorlinks=true,
  linkcolor=black,
  citecolor=black,
  urlcolor=blue!55!black,
  pdfauthor={Brian Tenneson},
  pdftitle={AMLD Security 3: Eight Formal Application Sketches}
}
\setlength{\headheight}{14pt}
\setlength{\parindent}{0pt}
\setlength{\parskip}{0.62em}
\setstretch{1.035}
\raggedbottom
\pagestyle{fancy}
\fancyhf{}
\fancyhead[L]{AMLD Security 3}
\fancyhead[R]{Companion research manuscript v1.0}
\fancyfoot[C]{\thepage}
\setcounter{tocdepth}{2}

\newtheorem{theorem}{Theorem}[chapter]
\newtheorem{proposition}[theorem]{Proposition}
\newtheorem{lemma}[theorem]{Lemma}
\newtheorem{corollary}[theorem]{Corollary}
\newtheorem{conjecture}[theorem]{Conjecture}
\theoremstyle{definition}
\newtheorem{definition}[theorem]{Definition}
\newtheorem{example}[theorem]{Example}
\theoremstyle{remark}
\newtheorem{remark}[theorem]{Remark}
\newtheorem{boundary}[theorem]{Boundary note}

\newcommand{\BANK}{\mathsf{BANK}}
\newcommand{\FB}{\mathsf{FUTUREBANK}}
\newcommand{\COMPASS}{\mathsf{COMPASS}}
\newcommand{\OCEAN}{\mathcal O}
\newcommand{\Req}{\operatorname{Req}}
\newcommand{\Impl}{\operatorname{Impl}}
\newcommand{\Reach}{\operatorname{Reach}}
\newcommand{\Acc}{\operatorname{Acc}}
\newcommand{\Ver}{\operatorname{Ver}}
\newcommand{\Cost}{\operatorname{Cost}}
\newcommand{\Risk}{\operatorname{Risk}}
\newcommand{\Dep}{\operatorname{Dep}}
\newcommand{\Future}{\operatorname{Future}}
\newcommand{\UNKNOWN}{\mathsf{UNKNOWN}}
\newcommand{\settle}{\mathsf{Settle}}
\newcommand{\Promote}{\mathsf{Promote}}
\newcommand{\Fire}{\mathsf{Fire}}

\begin{document}

\frontmatter
\hypersetup{pageanchor=false}
\begin{titlepage}
\centering
\vspace*{1.05in}
{\Huge\bfseries AMLD Security 3\par}
\vspace{0.25in}
{\Large Eight Formal Application Sketches for\par}
\vspace{0.08in}
{\Large Intelligent Target Settlement\par}
\vspace{0.45in}
{\large A complementary companion to\par}
\vspace{0.08in}
{\large\itshape Security by Separation, Verification, and Settlement\par}
\vspace{0.65in}
{\large Brian Tenneson\par}
\vspace{0.22in}
{\large Companion research manuscript v1.0\par}
\vspace{0.08in}
{\large September 2, 2026\par}
\vfill
\begin{minipage}{0.84\textwidth}
\small
The older treatise develops the defensive foundations: separation, identity, threat models,
containment, formal verification, the Ocean gate, and verifier-governed AMLD memory. This
companion begins where that book deliberately stops. It asks how the same architecture could be
applied, cautiously and formally, when an intelligent target seeker is aimed at cryptographic,
software, scientific, engineering, infrastructure, or AI-safety questions. Each chapter moves from
motivation to an explicit mathematical sketch, a modest proposition or conjecture, and a boundary
statement. None of the sketches establishes that a deployed system is vulnerable or that AMLD has
already achieved the proposed capability.
\end{minipage}
\end{titlepage}
\hypersetup{pageanchor=true}

\chapter*{Companion Statement}
\addcontentsline{toc}{chapter}{Companion Statement}
This manuscript is complementary in a precise sense. The older sibling, \emph{Security by
Separation, Verification, and Settlement} \cite{TennesonSecurity}, moves from elementary security
concepts toward AMLD. The present manuscript assumes those concepts and moves outward from AMLD
toward eight possible application domains. The two books therefore answer different questions:

\begin{center}
\begin{tabularx}{0.93\textwidth}{>{\raggedright\arraybackslash}p{0.25\textwidth}X X}
\toprule
 & Older sibling & This companion \\
\midrule
Primary question & What defensive architecture should exist before powerful reasoning is trusted? &
What could verifier-gated target settlement plausibly do in specific domains? \\
Direction & Security foundations $\rightarrow$ AMLD & AMLD $\rightarrow$ application sketches \\
Main result type & Defensive theorem & Conditional application theorem or research conjecture \\
Recurring warning & A verifier is only as good as its model and specification. & A formal pattern is not evidence of an achieved capability. \\
\bottomrule
\end{tabularx}
\end{center}

The shared writing rule remains \emph{motivation first, formalism immediately after}. The shared
security policy remains defensive: the mathematics is organized around auditing, containment,
certification, repair, and controlled deployment, not around operational exploitation.

\chapter*{Abstract}
\addcontentsline{toc}{chapter}{Abstract}
An automated target prover, theorem prover, or metalogical decider becomes security-relevant when
its target is a vulnerability, an invariant, a compressed state representation, or a certificate about
a protected system. The central problem is not merely computational power. It is how to separate
creative search from the authority to promote claims into trusted state. This manuscript gives a
common verifier-gated formalism and applies it to eight domains: cryptography, vulnerability
research, critical infrastructure, military and intelligence compartmentation, software correctness,
scientific discovery, engineering optimization, and AI safety. Across the applications, four motifs
recur: certificate-level settlement by the $P/R/I/C$ roles; separation of verified \(\BANK\) state from
speculative \(\FB\) state; future-sufficient quotienting; and Ocean-style regression gates for
candidate machinery. The results are intentionally conditional. They identify mathematical forms
that a defensible application could take and expose the preservation obligations that would have to
be proved. They do not report a new cryptanalytic attack, a universal vulnerability finder, or a proof
of global AI safety.

\tableofcontents

\mainmatter
\part{The Shared Formal Spine}

\chapter{From Intelligent Target Seeking to Verified Settlement}
\label{ch:spine}

\section{Targets, certificates, and environments}
Let a target instance be
\[
\mathfrak I=(c,\Gamma,\mathcal E,V,\mathcal B),
\]
where \(c\) is a claim or goal, \(\Gamma\) is the admitted theory or specification,
\(\mathcal E\) is the declared environment, \(V\) is an independent verifier, and
\(\mathcal B\) is a resource budget. Search may be learned, reflective, randomized, federated, or
otherwise adaptive. Validity is narrower:
\[
\operatorname{Accept}(\pi,c\mid\Gamma,\mathcal E)
\iff V(\pi,c,\Gamma,\mathcal E)=1.
\]
Changing a search policy may change which certificate is found and how quickly it is found; it must
not silently change the acceptance relation.

\section{Five operational outcomes}
For a formal target \(c\), an AMLD-style process distinguishes
\[
\settle(c)\in\{P,R,I,C,\UNKNOWN\}.
\]
Here \(P\) supplies a verified certificate for \(c\), \(R\) supplies one for \(\neg c\), \(I\)
supplies an independence certificate in a declared metatheory, and \(C\) supplies an inconsistency
certificate for the admitted assumptions. The fifth outcome records resource-bounded nonsettlement.
It is crucial that
\[
\UNKNOWN\not\Rightarrow c,
\qquad
\UNKNOWN\not\Rightarrow\neg c,
\qquad
\UNKNOWN\not\Rightarrow I(c).
\]

\section{Epistemic separation}
Let \(B_t\) be verified shared state and \(F_t\) speculative state at time \(t\). Candidate objects
may be generated freely into \(F_t\), but promotion is restricted:
\[
B_{t+1}=B_t\cup V_t(A_t),
\qquad
V_t(A_t):=\{a\in A_t:V(a,\Gamma_t,\mathcal E_t)=1\}.
\]
This is the mathematical core of \(\BANK/\FB\) separation. Provenance is retained, but provenance
does not substitute for verification.

\begin{proposition}[Verifier-gated noncontamination]
\label{prop:noncontamination}
Suppose the displayed update is the only rule that adds semantic assertions to \(B_t\). Then every
new assertion in \(B_{t+1}\setminus B_t\) has been accepted by \(V_t\).
\end{proposition}

\begin{proof}
By the update rule, \(B_{t+1}\setminus B_t\subseteq V_t(A_t)\). Membership in \(V_t(A_t)\)
is defined by verifier acceptance.
\end{proof}

\section{Future-sufficient quotienting}
Let \(\mathcal H\) be a set of search histories and \(Q:\mathcal H\to\mathcal S\) a descriptor.
The equivalence relation \(h_1\sim_Q h_2\) iff \(Q(h_1)=Q(h_2)\) is computationally useful only
if it preserves the future facts relevant to the target. One strong form is
\[
Q(h_1)=Q(h_2)
\Longrightarrow
\Future(h_1)\equiv_V\Future(h_2),
\]
where \(\equiv_V\) means equality of verifier-relevant continuations, dependency obligations,
scope, and certificate-lifting behavior.

\begin{proposition}[Sound quotient reuse]
\label{prop:quotient-reuse}
Assume each quotient certificate \(\bar\pi\) accepted at state \(Q(h)\) has a refinement map
\(L_h\) with
\[
\bar V(\bar\pi,Q(h))=1\Longrightarrow V(L_h\bar\pi,h)=1.
\]
Then reusing accepted quotient computation cannot introduce a false promoted claim in the original
state space.
\end{proposition}

\begin{proof}
Every promotion from a quotient state is refined by \(L_h\) and checked by \(V\). Soundness is
therefore inherited from the stated lifting condition, not from the quotient heuristic itself.
\end{proof}

\section{Guidance and the Ocean gate}
A \(\COMPASS\)-type score may prioritize actions,
\[
S(a\mid h)=w^{\mathsf T}\phi(a,h)-\lambda\widehat{\Cost}(a\mid h)
-\mu\widehat{\Risk}(a\mid h),
\]
but \(S\) is not a truth predicate. Candidate search machinery \(M\) can itself be tested against a
versioned Ocean suite \(\OCEAN_t\). If \(I_1,\ldots,I_k\) are hard invariants, a conservative gate is
\[
\Acc(M,\OCEAN_t)=1
\iff
\bigwedge_{o\in\OCEAN_t}\bigwedge_{j=1}^{k}I_j(M,o)=1.
\]
A verified failure \(r_t\) may then be retained by \(\OCEAN_{t+1}=\OCEAN_t\cup\{r_t\}\).
Passing a finite gate excludes tested failures; it does not prove universal safety.

\chapter{Application Map}

\begin{longtable}{>{\raggedright\arraybackslash}p{0.15\textwidth}>{\raggedright\arraybackslash}p{0.22\textwidth}>{\raggedright\arraybackslash}p{0.25\textwidth}>{\raggedright\arraybackslash}p{0.25\textwidth}}
\toprule
Domain & Formal target & Useful certificate & Central preservation obligation \\
\midrule
\endfirsthead
\toprule
Domain & Formal target & Useful certificate & Central preservation obligation \\
\midrule
\endhead
Cryptography & Hardness-relevant invariant or quotient & Proof that a reduction or lifting map preserves the target problem & Do not merge histories with different verifier-relevant futures \\
Vulnerability research & Conformance of implementation to policy & Checked violating trace or inductive proof & Model the complete sufficient routes and system state \\
Critical infrastructure & Bounded compromise propagation & Reachability certificate and non-regression record & Containment must not destroy required availability \\
Secure domains & Authorized boundary crossing and noninterference & Invariant proof or violating trace & Preserve labels, authority, provenance, and declassification rules \\
Software correctness & Safety invariant over reachable states & Inductive proof or finite counterexample & Checker and dependency closure remain fixed \\
Scientific discovery & Predictively sufficient low-dimensional state & Theorem, checked derivation, or preregistered empirical certificate & Preserve the declared observable family, not merely training fit \\
Engineering optimization & Feasible and objective-preserving state merge & Correspondence between feasible continuations & Preserve downstream constraints and optimum values \\
AI safety & Separation of generation, verification, testing, and authority & Gate receipts and invariant certificates & Policy revision must not revise its own truth conditions silently \\
\bottomrule
\end{longtable}

The table is not evidence that all eight applications are equally mature. It identifies the object
that would need to be certified before speculative target seeking could be treated as defensive
knowledge.

\part{Eight Plausible Applications}

\chapter{Cryptography}
\label{ch:crypto}

Discovering unexpected structure in a problem believed to be hard could matter to cryptography.
The defensible research question is not whether AMLD can simply try keys faster. It is whether a
small, independently checkable invariant or quotient can replace a much larger search without
changing the cryptographic problem being solved.

\section{Plausible application}
Let \(X_\kappa\) be an instance family at security parameter \(\kappa\), let \(\mathcal H_\kappa\)
be legal search histories, and let \(T(h,a)\) be continuation by action \(a\). A quotient candidate
\(Q_\kappa:\mathcal H_\kappa\to\mathcal S_\kappa\) induces
\[
h_1\sim_Qh_2\iff Q_\kappa(h_1)=Q_\kappa(h_2).
\]
For defensive significance, equality of quotient states must imply more than similar local features.
It should preserve terminal status, admissible continuations, certificate dependencies, and success
probability under the declared model. A strong deterministic obligation is a verifier-respecting
bisimulation:
\[
h_1\sim_Qh_2\Longrightarrow
\left\{
\begin{array}{l}
V_{\rm term}(h_1)=V_{\rm term}(h_2),\\
\forall a_1\in A(h_1)\;\exists a_2\in A(h_2):T(h_1,a_1)\sim_QT(h_2,a_2),\\
\forall a_2\in A(h_2)\;\exists a_1\in A(h_1):T(h_1,a_1)\sim_QT(h_2,a_2).
\end{array}\right.
\]
If the search is stochastic, the corresponding condition is equality of the pushed-forward
transition kernels on \(\mathcal S_\kappa\). A \(\COMPASS\)-like policy may search over candidate
descriptors by a score
\[
J(Q)=\widehat{\operatorname{reuse}}(Q)-\lambda\log|\operatorname{im}Q|
-\mu\widehat{\Cost}(Q)-\nu\widehat{\operatorname{collision\ risk}}(Q),
\]
but promotion requires a certificate for the preservation obligation. Under that certificate, dynamic
programming or lemma reuse may occur once per quotient class rather than once per history. The
scientific object is therefore a proved structural compression of an admitted problem, not an
unverified claim that a named encryption scheme is weak.

\begin{conjecture}[Certified quotient advantage]
For some structured cryptographic audit families, there exists a future-sufficient \(Q\) for which
verified quotient search uses asymptotically fewer distinct continuation states than history-level
search while returning only liftable certificates.
\end{conjecture}

\begin{boundary}
No argument in this chapter establishes an efficient attack on AES, RSA, elliptic-curve cryptography,
or a post-quantum construction. Hardness is scheme-, model-, and parameter-specific.
\end{boundary}

\paragraph{Source anchors.} Safe quotienting and \(\COMPASS\) are developed in
\emph{What Checks the Proof?} v6.80, pp. 28--29 and 68 \cite{TennesonChecks}; security as a
settlement problem appears in the older security treatise, pp. 42--43 \cite{TennesonSecurity}.

\chapter{Cybersecurity and Vulnerability Research}
\label{ch:vulnerability}

A target-seeking system could generate many suspected traces. Defensively, the important move is to
replace the vague target ``find a vulnerability'' by a model-relative conformance claim and a
machine-checkable defect certificate. Speculation remains useful, but it remains quarantined.

\section{Plausible application}
Let \(\mathcal M=(S,S_0,A,\to)\) be a labeled transition system and let \(Z\) be its complete finite
or infinite traces. A policy defines \(\Req(z)\in\{0,1\}\); the implementation semantics defines
\(\Impl(z)\in\{0,1\}\). One-sided acceptance soundness is
\[
\Phi_{\rm safe}:=\forall z\in Z\,[\Impl(z)\Rightarrow\Req(z)].
\]
If a finite trace \(z^\star=(s_0,a_0,\ldots,s_n)\) carries receipts proving
\[
s_0\in S_0,\quad s_i\xrightarrow{a_i}s_{i+1},\quad
\Impl(z^\star)=1,\quad \Req(z^\star)=0,
\]
then \(z^\star\) is a finite \(R\)-certificate. Conversely, a \(P\)-certificate may consist of an
inductive invariant \(I\) satisfying
\[
S_0\subseteq I,\qquad
(s\in I\land s\to s')\Rightarrow s'\in I,
\qquad
s\in I\Rightarrow\neg\mathsf{Bad}(s).
\]
The \(I\)-role asks a different question: whether \(\Phi_{\rm safe}\) is independent of the admitted
audit theory in a declared metatheory. The \(C\)-role searches for inconsistent policy obligations.
Candidate traces, invariants, and repairs reside in \(F_t\); only checker-accepted objects enter \(B_t\).
For an evolving analyzer \(M_t\), Ocean admission can require preservation of four architectural
invariants: rejected candidates never enter \(B_t\); every derived record has a receipt; dependencies
are closed; and resource limits are enforced. This converts vulnerability research into an
evidence-producing audit process in which the reasoning system proposes and prioritizes, while a
separate checker determines what has actually been established.

\begin{proposition}[Counterexample sufficiency]
If the displayed receipts for \(z^\star\) verify, then \(\Phi_{\rm safe}\) is false in the declared
transition model and policy.
\end{proposition}

\begin{proof}
The verified trace instantiates the universal quantifier with an implemented acceptance that is not
required, falsifying the implication.
\end{proof}

\begin{boundary}
A checked defect is relative to a model and specification. An incomplete model can miss real behavior;
an inadequate policy can be implemented perfectly and still protect the wrong property.
\end{boundary}

\paragraph{Source anchors.} The authentication-defect certificate and four-role audit appear in
\emph{What Checks the Proof?} v6.80, pp. 54--55 \cite{TennesonChecks}; the Ocean moat appears in
the older treatise, pp. 32--33 \cite{TennesonSecurity}.

\chapter{Critical Infrastructure}
\label{ch:infrastructure}

For hospitals, power systems, water utilities, transportation, and communications, an initial cyber
failure can become a physical or availability failure. The natural defensive object is therefore not
one secure/insecure bit but a graph of propagation, authority, service, and recovery.

\section{Plausible application}
Model the system as a directed labeled graph \(G=(V,E,\ell,w)\), where \(\ell(e)\) records the type
of transition and \(w(v)\ge0\) is a declared consequence weight. For a compromised component
\(c\), define
\[
B_G(c)=\Reach_G(c),
\qquad
L_G(c)=\sum_{v\in B_G(c)}w(v).
\]
A segmentation design \(D\) produces \(G_D=(V,E_D)\) together with a set \(R_D\) of service
paths that must remain available. If \(E_{D_2}\subseteq E_{D_1}\), reachability monotonicity gives
\[
B_{G_{D_2}}(c)\subseteq B_{G_{D_1}}(c),
\qquad
L_{G_{D_2}}(c)\le L_{G_{D_1}}(c).
\]
But removal of every edge is not a useful defense. Let
\[
\mathbf J(D)=\left(
p_{\rm comp}(D),
\max_cL_{G_D}(c),
A_{\rm loss}(D),
C_{\rm verify}(D),
C_{\rm repair}(D)
\right).
\]
A candidate dominates another only under a declared partial order and hard service constraints,
for example
\[
D_2\succ_XD_1
\iff
m_X(D_2)<m_X(D_1)
\land
\forall j\in J_{\rm protected},\;m_j(D_2)\le m_j(D_1)
\land R_{D_2}\models\Psi_{\rm service}.
\]
An AMLD controller can let \(P\) prove containment and service obligations, \(R\) find checked
violating paths, \(I\) test whether the obligations are independent of the model, and \(C\) detect
incompatible requirements. The useful output is a certificate-bearing Pareto comparison, not a
single heuristic security score.

\begin{proposition}[Edge-removal containment]
If \(E'\subseteq E\), then for every \(c\in V\), \(\Reach_{(V,E')}(c)\subseteq\Reach_{(V,E)}(c)\).
\end{proposition}

\begin{proof}
Every path using only edges in \(E'\) is also a path using edges in \(E\).
\end{proof}

\begin{boundary}
The proposition proves containment in the graph model only. It does not validate the graph, estimate
compromise probability, or show that required services survive segmentation.
\end{boundary}

\paragraph{Source anchors.} Blast radius and narrow bridges are treated in the older security
treatise, pp. 10--11; multi-objective security and separation-aware control appear on pp. 53--56
\cite{TennesonSecurity}.

\chapter{Military and Intelligence Compartmentation}
\label{ch:compartmentation}

Compartmentation gives a mathematically clean example of the book's security philosophy: separate
domains, restrict crossings, preserve provenance, and never grant speculative reasoning direct
operational authority. The formulation below is about defensive authorization and noninterference,
not operational attack planning.

\section{Plausible application}
Let \((L,\preceq)\) be a security-label lattice and let each state component \(x\) carry label
\(\lambda(x)\in L\). Partition the state space into domains
\[
V=\bigsqcup_{i=1}^{k}V_i.
\]
For every cross-domain transition \(e=(x,y)\), define a boundary predicate
\(P_e(s,\alpha,\rho)\), where \(s\) is current state, \(\alpha\) is asserted authority, and \(\rho\)
is provenance. The safety target is
\[
\Phi_{\rm cross}:=
\Box\left(
\Fire(e,s,\alpha,\rho)
\Rightarrow P_e(s,\alpha,\rho)
\right).
\]
For confidentiality, let \(s\equiv_\ell s'\) mean that states agree on observations visible at low
label \(\ell\). A possibilistic noninterference target is
\[
s_0\equiv_\ell s'_0
\Longrightarrow
\operatorname{Obs}_\ell(\operatorname{Run}(s_0,u))
=\operatorname{Obs}_\ell(\operatorname{Run}(s'_0,u))
\]
for every admitted low action sequence \(u\), except at explicitly modeled declassification gates.
Role \(P\) may prove these invariants compositionally; \(R\) may supply a checked boundary-crossing
or observational counterexample; \(I\) and \(C\) retain their metatheoretic meanings. Search output
does not become trusted merely because a designated role produced it:
\[
c\in B_{t+1}\setminus B_t
\Longrightarrow
\exists\pi\;V(\pi,c,\Gamma,\mathcal E)=1.
\]
Thus compartmentation is represented simultaneously as a graph boundary, an authorization
predicate, and an observational equivalence, with the verifier mediating any promotion into trusted
state.

\begin{conjecture}[Cross-domain shared-certificate advantage]
For some compartmented systems, a provenance-preserving shared \(\BANK\) of verified interface
lemmas reduces total audit cost relative to isolated audits without weakening domain-local acceptance.
\end{conjecture}

\begin{boundary}
Noninterference depends on the observation model, declassification rules, scheduler, and environment.
The displayed formula is not a universal characterization of national-security systems.
\end{boundary}

\paragraph{Source anchors.} Separation, least privilege, and narrow bridges are treated in the older
security treatise, pp. 9--11; $P/R/I/C$ security settlement appears on pp. 42--43
\cite{TennesonSecurity}.

\chapter{Software Correctness}
\label{ch:software}

Software correctness is the most direct beneficial application because both proof and refutation can
often be represented by finite machine-checkable objects. Learned guidance may choose what to try;
it need not be trusted to decide what counts as a valid proof.

\section{Plausible application}
Let \(\mathcal M=(S,S_0,\to)\) be a transition system and let \(I:S\to\{0,1\}\) be a safety
invariant. An inductive certificate proves
\[
\forall s\in S_0\;I(s),
\qquad
\forall s,s'\in S\;[I(s)\land s\to s']\Rightarrow I(s').
\]
It follows that \(I\) holds throughout \(\Reach_{\mathcal M}(S_0)\). Role \(P\) searches for such an
invariant, auxiliary lemmas, and checker receipts. Role \(R\) searches for a finite trace
\[
s_0\to s_1\to\cdots\to s_n,
\qquad s_0\in S_0,
\qquad \neg I(s_n).
\]
The reasoning architecture is also software and should satisfy its own executable invariants:
\[
u\in B_{t+1}\setminus B_t
\Rightarrow\exists r\;V(r,u)=1,
\]
\[
V(r,u)=0\Rightarrow u\notin B_{t+1},
\qquad
u\in B_t\Rightarrow\Dep(u)\subseteq B_t.
\]
If verifier receipts are content-addressed and dependencies are immutable, replay can check that
the accepted object still means what it meant at promotion time. Candidate proofs and candidate
program repairs remain in \(F_t\) until checking succeeds. An Ocean suite then includes historical
counterexamples, malformed certificates, resource-limit tests, and dependency-closure cases. This
creates two mutually reinforcing proof obligations: the application program must preserve its safety
invariant, and the settlement machinery must preserve the epistemic boundary by which that claim is
accepted.

\begin{proposition}[Inductive reachability safety]
If the base and step obligations above hold, then every reachable state satisfies \(I\).
\end{proposition}

\begin{proof}
Induct on the length of a path from \(S_0\). The base case uses \(S_0\subseteq I\); the step uses
transition preservation.
\end{proof}

\begin{boundary}
The theorem establishes only the invariant encoded by \(I\) in the modeled semantics. It does not
prove that the model matches compiled code, hardware, operators, or every environmental condition.
\end{boundary}

\paragraph{Source anchors.} Policy-neutral verification and executable architecture invariants appear
in \emph{What Checks the Proof?} v6.80, pp. 10 and 69--70 \cite{TennesonChecks}.

\chapter{Scientific Discovery}
\label{ch:science}

The Quotient Hunter idea is scientifically interesting when many apparent histories can be replaced
by a smaller state that preserves everything relevant to future prediction or verification. Invariant
discovery then becomes a search for sufficient representation, not merely a search for correlation.

\section{Plausible application}
Let \(H_t\) denote an observed history and \(Y_{t:\infty}\) a declared family of future observables.
A descriptor \(Q(H_t)\) is predictively sufficient when
\[
\mathcal L(Y_{t:\infty}\mid H_t)
=\mathcal L(Y_{t:\infty}\mid Q(H_t)).
\]
Equivalently, under regularity conditions,
\[
Y_{t:\infty}\perp H_t\mid Q(H_t).
\]
For deterministic models and observable family \(\mathcal F\), the obligation may instead be
\[
Q(h_1)=Q(h_2)
\Rightarrow
\forall f\in\mathcal F\;\forall\tau\ge0,
\quad f(Y_{h_1}(\tau))=f(Y_{h_2}(\tau)).
\]
Candidate descriptors can be ranked by an information-compression objective
\[
U(Q)=\widehat I(Q(H_t);Y_{t:\infty})
-\lambda I(Q(H_t);H_t)
-\mu\widehat{\Cost}(Q),
\]
or by an exact symbolic compression metric when a formal model is available. Yet \(U(Q)\) is only
a search score. A descriptor enters verified state only after a domain-appropriate certificate: a
proof in a formal model, an exact algebraic derivation, or a preregistered empirical test with a clearly
limited conclusion. \(\FB\) may retain failed invariants and conditional ``what if'' models without
asserting them. If a descriptor passes the preservation gate, downstream inference may reuse the same
reasoning for every history in a quotient class. The payoff is a certified reduction of effective state
dimension and an explicit statement of which observables survive the reduction.

\begin{conjecture}[Verifier-history discovery advantage]
On some repeated scientific-model families, retaining rejected quotient hypotheses with provenance
reduces expected rediscovery cost without increasing false certification, relative to retaining only
accepted descriptors.
\end{conjecture}

\begin{boundary}
Predictive sufficiency is relative to the chosen observables and data-generating assumptions. Good
holdout prediction alone is not a theorem that the quotient is universally correct.
\end{boundary}

\paragraph{Source anchors.} \(\FB\), planning, and \(\COMPASS\) appear in
\emph{What Checks the Proof?} v6.80, pp. 27--30; future-sufficient quotienting appears on p. 68
\cite{TennesonChecks}.

\chapter{Engineering Optimization}
\label{ch:engineering}

Scheduling, placement, routing, circuit design, materials design, and configuration all contain large
combinatorial spaces. Safe state merging can reduce repeated work only when it preserves future
constraints and objective values.

\section{Plausible application}
Consider
\[
\min_{x\in X} f(x)
\quad\text{subject to}\quad
g_i(x)\le0\;(i=1,\ldots,m),
\qquad h_j(x)=0\;(j=1,\ldots,r).
\]
For a partial design history \(h\), let
\[
\mathcal C(h)=\{x\in X:x\text{ extends }h\text{ and satisfies inherited constraints}\}.
\]
A strong exact equivalence is
\[
h_1\sim h_2
\iff
\exists\varphi:\mathcal C(h_1)\to\mathcal C(h_2)\text{ bijective such that }
f(x)=f(\varphi(x))
\]
and all checker-relevant labels are preserved. A weaker quotient may preserve only feasibility and an
optimal value:
\[
Q(h_1)=Q(h_2)
\Rightarrow
\left[
\mathcal C(h_1)=\varnothing\iff\mathcal C(h_2)=\varnothing
\right]
\land
\inf_{x\in\mathcal C(h_1)}f(x)=\inf_{x\in\mathcal C(h_2)}f(x).
\]
Under either certified obligation, dynamic programming may cache downstream value and feasibility
by quotient state. A \(\COMPASS\)-like scheduler may prioritize
\[
S(a\mid h)=
\alpha\Delta_{\rm goal}+\beta\rho_Q+\gamma\operatorname{reuse}
-\delta\widehat{\Cost}-\eta\widehat{\Risk},
\]
while feasibility remains checker-defined:
\[
\operatorname{Feas}(x)
\iff
\bigwedge_i g_i(x)\le0
\land\bigwedge_j h_j(x)=0
\land V(x,\Gamma)=1.
\]
Learned guidance therefore decides where to spend computation; it does not redefine admissibility.
The research question is whether certified equivalence classes can be discovered cheaply enough that
their reuse benefit exceeds the cost of proving them.

\begin{proposition}[Optimal-value reuse]
If \(Q\) satisfies the displayed optimal-value preservation obligation, then an exact cached value at
\(Q(h)\) is valid for every history in that quotient class.
\end{proposition}

\begin{proof}
The obligation states equality of the infimum for any two representatives of the same class.
\end{proof}

\begin{boundary}
Preserving an optimum value does not necessarily preserve an optimizer, robustness margin, or
implementation cost. Those quantities must be added to the equivalence obligation when relevant.
\end{boundary}

\paragraph{Source anchors.} \(\COMPASS\) and safe quotienting appear in
\emph{What Checks the Proof?} v6.80, pp. 28--29 and 68 \cite{TennesonChecks}.

\chapter{AI Safety Itself}
\label{ch:ai}

A shortcut-generating system is materially different from a shortcut-generating system that is also
allowed to declare its own outputs valid. The application here is architectural: separate generation,
testing, verification, memory promotion, and operational authority.

\section{Plausible application}
Let \(G_\theta\) be a mutable generator or search policy and \(V_\omega\) a separately governed
verifier. For a claim \(c\), certificate \(\pi\), theory \(\Gamma\), and environment \(\mathcal E\),
define
\[
\Promote(c,\pi)
\iff V_\omega(c,\pi,\Gamma,\mathcal E)=1.
\]
Policy-neutral verification requires generator changes not to alter acceptance:
\[
\forall\theta_1,\theta_2,c,\pi,
\quad
V_\omega(c,\pi,\Gamma,\mathcal E)
\text{ is independent of whether }G_{\theta_1}\text{ or }G_{\theta_2}\text{ proposed }(c,\pi).
\]
Candidate self-modifications \(\theta'\) remain speculative,
\[
\theta'\in F_t\not\Rightarrow\theta'\in B_t,
\]
and must pass hard invariants \(H_1,\ldots,H_k\) plus an Ocean regression suite:
\[
\operatorname{Gate}(\theta')=1
\iff
\left(\bigwedge_{o\in\OCEAN_t}\bigwedge_{j=1}^{k}H_j(G_{\theta'},o)=1\right)
\land V_\omega(\pi_{\theta'},\theta',\Gamma_{\rm deploy})=1.
\]
The gate may also include a statistical endpoint, but statistical success cannot override a failed hard
invariant. A verified regression failure \(r_t\) updates \(\OCEAN_{t+1}=\OCEAN_t\cup\{r_t\}\).
Operational authority is granted by a distinct deployment controller with least privilege and rollback,
not by membership in a speculative search frontier. This does not solve alignment in general. It
supplies a checkable separation law: policy may change attention, but promotion remains conditional
on an independently governed acceptance relation.

\begin{proposition}[Generator-independent promotion]
Under the displayed promotion rule, changing \(\theta\) alone cannot cause a certificate rejected by
\(V_\omega\) to enter \(B_t\) through the modeled update.
\end{proposition}

\begin{proof}
Promotion depends on \(V_\omega\), not on the identity or parameters of the proposing generator.
Apply Proposition~\ref{prop:noncontamination}.
\end{proof}

\begin{boundary}
Independent verification is a safety mechanism, not a proof of global AI safety. The verifier,
specification, interfaces, deployment controller, and environment remain possible failure surfaces.
\end{boundary}

\paragraph{Source anchors.} Policy-neutral verification and safe quotienting appear in
\emph{What Checks the Proof?} v6.80, pp. 10 and 68 \cite{TennesonChecks}; the Ocean moat appears
in the older security treatise, pp. 32--33 \cite{TennesonSecurity}.

\part{Research Boundaries and Program}

\chapter{A Cross-Domain Experimental Program}

The eight sketches can be tested without pretending that they are already established capabilities.
For each application family, freeze
\[
(\text{target family},\text{model},\text{verifier},\text{budget},\text{baseline},
\text{split},\text{stopping rule},\text{metric}).
\]
Compare at least four conditions: unguided certified search; guided search without quotient reuse;
guided search with candidate quotienting but no promotion; and guided search with verifier-certified
quotient reuse. A useful primary endpoint is verified settlement per unit resource. Secondary
endpoints include certificate size, wall time, peak memory, repeated-search cost, false-promotion
count, and non-regression coordinates such as availability.

\section{A minimum preregistration schema}
\begin{enumerate}[leftmargin=*]
\item State the target claim and the exact verifier acceptance relation.
\item State which data or instances inform search and which remain held out.
\item State the quotient preservation obligation before observing target results.
\item Match baselines on verifier, budget, and stopping conditions.
\item Record every promoted object with its certificate, dependencies, provenance, and environment.
\item Treat \(\UNKNOWN\) as a distinct outcome rather than silently reclassifying it.
\item Retain verified failures as regression cases and report all non-regression coordinates.
\end{enumerate}

\begin{conjecture}[Separation-aware advantage]
On some repeated target-settlement distributions, a controller that optimizes both settlement progress
and verifier-boundary preservation attains a better Pareto frontier than a matched controller that
optimizes settlement success alone.
\end{conjecture}

\chapter{Claims This Manuscript Does Not Make}

\begin{enumerate}[leftmargin=*]
\item No finite Ocean suite proves universal security.
\item No heuristic score is promoted to a truth predicate.
\item No failure to settle is treated as proof of safety, falsity, or independence.
\item No quotient is safe merely because it compresses a benchmark well.
\item No proof of implementation conformance establishes adequacy of the specification.
\item No component-local proof automatically establishes secure composition.
\item No chapter reports a new practical cryptanalytic attack or operational vulnerability.
\item No verifier-gated architecture, by itself, proves global AI safety.
\end{enumerate}

The positive claim is narrower and, for research, more useful: each application can be formulated so
that the powerful component proposes search directions while an explicit preservation obligation and
independent checker govern promotion. That is the precise point at which this companion joins its
older sibling.

\appendix
\chapter{Glossary}

\begin{description}[style=nextline,leftmargin=1.8em,labelindent=0pt]
\item[AMLD] A verifier-gated, multi-role architecture for settling formal targets through proof,
refutation, independence, contradiction, or an explicit resource-bounded unknown outcome.
\item[Blast radius] The set or weighted consequence of resources reachable after a modeled compromise.
\item[Certificate] A finite object whose validity for a declared claim can be checked independently.
\item[COMPASS] Learned or engineered guidance that ranks search actions without defining truth.
\item[Future-sufficient quotient] A state compression that identifies histories only when all declared
future facts needed for search and verification are preserved.
\item[Ocean] A versioned adversarial and regression environment used before candidate machinery is
granted trusted authority.
\item[Policy-neutral verification] The requirement that changing search policy does not silently change
the acceptance relation.
\item[Settlement] Acceptance of a certificate for one of the declared $P/R/I/C$ outcomes.
\item[Verified BANK] Shared trusted state whose semantic assertions enter only through the declared
verification gate.
\item[FUTUREBANK] Quarantined speculative state containing hypotheses, imagined traces, candidate
quotients, and proposed repairs not yet admitted to the verified BANK.
\end{description}

\chapter{Symbols and Notation}

\begin{longtable}{>{\raggedright\arraybackslash}p{0.18\textwidth}>{\raggedright\arraybackslash}p{0.68\textwidth}}
\toprule
Symbol & Meaning in this manuscript \\
\midrule
\endfirsthead
\toprule
Symbol & Meaning in this manuscript \\
\midrule
\endhead
\(c\) & Target claim or goal \\
\(\Gamma\) & Admitted theory, specification, and scoped dependencies \\
\(\mathcal E\) & Declared environment or threat model \\
\(V\) & Independent verifier or checker \\
\(P,R,I,C\) & Prover, Refuter, Independence, and Contradiction settlement roles \\
\(\UNKNOWN\) & No accepted settlement certificate within declared conditions \\
\(B_t\) & Verified BANK contents at time \(t\) \\
\(F_t\) & FUTUREBANK speculative contents at time \(t\) \\
\(Q\) & Candidate future-sufficient descriptor or quotient map \\
\(h_1\sim_Qh_2\) & Histories have the same quotient state \\
\(\OCEAN_t\) & Ocean regression suite at time \(t\) \\
\(\Reach_G(c)\) & Vertices reachable from \(c\) in graph \(G\) \\
\(D_2\succ_XD_1\) & \(D_2\) strictly improves the declared metric for vulnerability class \(X\), subject to stated non-regression constraints \\
\bottomrule
\end{longtable}

\backmatter
\chapter*{References}
\addcontentsline{toc}{chapter}{References}
\begin{thebibliography}{99}

\bibitem{Anderson}
Ross Anderson.
\newblock \emph{Security Engineering: A Guide to Building Dependable Distributed Systems}.
\newblock 3rd ed., Wiley, 2020.

\bibitem{Bishop}
Matt Bishop.
\newblock \emph{Computer Security: Art and Science}.
\newblock 2nd ed., Addison-Wesley, 2018.

\bibitem{CoverThomas}
Thomas M. Cover and Joy A. Thomas.
\newblock \emph{Elements of Information Theory}.
\newblock 2nd ed., Wiley, 2006.

\bibitem{SaltzerSchroeder}
Jerome H. Saltzer and Michael D. Schroeder.
\newblock ``The Protection of Information in Computer Systems.''
\newblock \emph{Proceedings of the IEEE} 63(9), 1975, 1278--1308.

\bibitem{TennesonSecurity}
Brian Tenneson.
\newblock \href{https://btenneson.github.io/pub/papers/amld_security_treatise_v0_1/}
{\emph{Security by Separation, Verification, and Settlement: An AMLD-Oriented Mathematical Treatise}}.
\newblock Working research manuscript v0.1, August 30, 2026.

\bibitem{TennesonChecks}
Brian Tenneson.
\newblock \href{https://btenneson.github.io/pub/papers/what_checks_the_proof_v6_80/}
{\emph{What Checks the Proof? A Textbook Introduction to Automated Mathematical Logic Deciders}}.
\newblock Version 6.80, August 22, 2026.

\end{thebibliography}

\clearpage
\thispagestyle{empty}
\vspace*{0.66\textheight}
\begin{center}
\fbox{\begin{minipage}{0.82\textwidth}
\itshape\small\raggedright
AI assistance note. This manuscript develops Brian Tenneson's AMLD, BANK/FUTUREBANK,
COMPASS, Quotient Hunter, Ocean, and verifier-gated settlement ideas. ChatGPT assisted with
organization, mathematical formalization, LaTeX drafting, and editorial revision. The named author
is responsible for reviewing the definitions, propositions, conjectures, citations, security boundaries,
and final publication text.
\end{minipage}}
\end{center}

\end{document}
